\section{\method Preface: Harvesting Jailbreak Tactics In-the-Wild}

Given a target model $M_{\text{target}}$, a harmful prompt $\mathcal{P}$ will elicit either a harmful ($\mathcal{P}\mathcal{R}_{h}^{M_{\text{target}}}$) or a benign ($\mathcal{P}\mathcal{R}_{b}^{M_{\text{target}}}$) model response. The goal of red-teaming is to identify harmful prompts $\mathcal{P}$ that reveal harmful responses from $M_{\text{target}}$. Jailbreaking, a more challenging form of red-teaming, aims to revise known harmful prompts $\mathcal{P}$ that currently elicit benign responses into \textit{adversarial} prompts $\mathcal{AP}$ to 
bypass the safeguard of $M_{\text{target}}$ by eliciting target harmful responses instead.

Our current knowledge of \textit{jailbreak tactics} used in forming adversarial attacks is relatively limited, and recent works uncover a narrow range of possible jailbreaks \citep{pap,pair,tap,rainbowteaming}. 
To overcome this limitation, we mine real-world chat logs, which is a surprisingly rich source of diverse jailbreak tactics, even though these users were not specifically instructed to jailbreak the system.

\begin{figure}[t]
    \captionof{table}{(Left) shows the number of items (\textbf{Total}), number of deduplicated unique clusters (\textbf{Uniq.}), and per query count (\textbf{Per.}) for jailbreak tactics automatically mined from \itw user queries in \lmsysshort and \wildchatshort, which contain a greater diversity and quantity of jailbreak tactics compared to those from other sources. \underline{Underline} indicates a sub-sampled set of queries. (Right) shows the top common jailbreak tactics and their percentage of occurrence.}
    \label{tab:compare_strategy_diversity}
    
  \begin{minipage}{.65\linewidth}
    \centering
    \small
    
    \begin{tabular}{
    l@{\hspace{0.6\tabcolsep}}
    l@{\hspace{0.05\tabcolsep}}
    r@{\hspace{0.95\tabcolsep}}
    r@{\hspace{0.95\tabcolsep}}
    r@{\hspace{0.95\tabcolsep}}
    r@{\hspace{0.95\tabcolsep}}
    }
        \toprule
        
        \multicolumn{2}{c}{\textbf{Data Source}} & \multicolumn{1}{c}{\textbf{Query}} &  \multicolumn{3}{c}{\textbf{Jailbreak Tactics}} \\
        \cmidrule(r){1-2} \cmidrule(r){3-3} \cmidrule(r){4-6}

        \textbf{Type} & \textbf{Name} & \textbf{Total} & \textbf{Total} & \textbf{Uniq.} &  \textbf{Per.} \\
        \midrule
        
        \multirow{3}{*}{\makecell[tl]{\itwshort}} & \lmsysshort \citep{zheng2023lmsyschat1m} & 7,873 & 43,220 & 2,526 & 5.49 \\
        & \wildchatshort \citep{zhao2024inthewildchat} & 8,981 & 62,218 & 3,903 & 6.93\\

        & \cellcolor{columbiablue} Combined & \cellcolor{columbiablue} 16,854 & \cellcolor{columbiablue} 105,438 & \cellcolor{columbiablue} 5,688 & \cellcolor{columbiablue} 6.26 \\
        \midrule

        \multirow{3}{*}{\makecell[tl]{Jailbreak \\ Templates}} & \danshort \citep{doanytingnow} & 666 & 4,378 & 510 & 6.57 \\
        & \trustllm \citep{sun2024trustllm} & 1,400 & 4,531 & 280 & 3.24 \\        
        & \decodingtrust \citep{wang2023decodingtrust} & 5 & 8 & 5 & 1.60 \\        
        \midrule  
        
        \multirow{3}{*}{\makecell[tl]{Semantic \\ Jailbreak \\ Methods}} & PAIR \citep{pair} & \underline{400} & 1,854 & 162 & 4.64 \\
        & TAP \citep{tap} & \underline{398} & 1,861 & 149 & 4.68 \\       
        & PAP \citep{pap} & \underline{398} & 1,564 & 118 & 3.93 \\
        \midrule 

        \multirow{3}{*}{\makecell[tl]{Safety \\ Training \\ Data}} & \hhrlhf \citep{hh-rlhf} & \underline{500} & 884 & 66 & 1.77 \\
        & \safellamashort \citep{safetytuned_llama} & \underline{500} & 911 & 66 & 1.82 \\   
        & Safe-RLHF \citep{dai2023saferlhf} & \underline{500} & 1,034 & 84 & 2.07 \\   
        \bottomrule

    \end{tabular}
    % \captionof{table}{Table caption}
    \end{minipage}\hfill
    \begin{minipage}{0.3\linewidth}
    \includegraphics[width=1\linewidth]{figures/top_tactics.pdf}
  \end{minipage}

\end{figure}


\subsection{Mining Jailbreak Tactics from Real-World User-Chatbot Interactions}

With a seed set of manually-identified tactics, we apply \gptfour to expand the discovery automatically.

\textbf{Gathering \itwshort User-Written Adversarial Harmful Prompts.} 
We first collect candidate adversarial prompts from all single-turn conversations in \lmsysshort \citep{zheng2023lmsyschat1m} and \wildchatshort \citep{zhao2024inthewildchat} that are flagged by the OpenAI Moderation API.\footnote{We include conversations where either the user prompt or the model response is flagged as harmful, as the OpenAI moderation API sometimes fails to flag nuanced and hard-to-detect unsafe user prompts.} We then filter out trivial non-adversarial prompts by feeding candidates through a lightly safety-trained model (\tulu-7B), keeping those that elicit harmful model responses as judged by the \llamaguard safety classifier \citep{llamaguard}; this yields 16,850 final prompts.

\textbf{Identifying Seed Jailbreak Tactics by Manual Examination.} 
We manually examine $\sim$200 \itwshort prompts sampled from our \itwshort adversarial prompt set to identify 35 seed jailbreak tactics with definitions (see the full list in Table \ref{tab:appendix_example_tactics} and \ref{tab:appendix_example_tactics_cont} in \S\ref{assec:manually_mined}).

\textbf{Automatic Tactics Discovery Aided by \gptfour.} 
With seed jailbreak tactics, we apply \gptfour to scale the tactic mining. For each adversarial prompt, \gptfour is given two tasks: (1) extracting the core vanilla request; (2) identifying both \textit{existing} and potentially \textit{novel} jailbreak tactics in the adversarial prompt. \gptfour additionally identifies an \textit{excerpt} corresponding to each tactic, the \textit{definition} to describe novel tactic, and \textit{reasoning} of why the tactic applies. 
Each step is carefully prompted with a demonstration example (see prompts in Table \ref{tab:strategy_mining_prompt} and \ref{tab:adversarial_prompt_simplification_prompt} in \S \ref{assec:tactics_mining_prompts}). We then deduplicate all tactics by clustering on their corresponding definitions\footnote{The deduplication is done on tactic definitions instead of names, as we observe that tactics with drastically different names may capture the same definition.} with sentence embeddings\footnote{Sentence embeddings are obtained from Nomic Embed (\url{https://huggingface.co/nomic-ai/nomic-embed-text-v1}) using clustering threshold of 0.75. Examples of tactic clusters are shown in Table \ref{tab:tactic_cluster_examples} in Appendix \ref{asec:strategy_mine}.} and report the statistics of these unique clusters in Table \ref{tab:compare_strategy_diversity}.

\subsection{What Tactics Are Adopted by In-the-Wild Users for Jailbreaking LLMs?}

Table \ref{tab:compare_strategy_diversity} shows the top \itw jailbreak tactics, including a mixture of stylistic, syntactic, formatting, writing genre, and context-based tricks.
Specifically, it uncovers novel tactics not systematically documented previously, such as ``prefacing the harmful content with a content warning or disclaimer,'' ``setting blame for non-compliance,'' or ``cloaking harm in humor'' (more examples of novel tactics in Table \ref{tab:appendix_example_auto_tactics} of Appendix \S \ref{assec:tactics_mining_prompts}).

In addition, as shown in Table \ref{tab:compare_strategy_diversity}, \itwshort adversarial user queries contain the richest set of unique jailbreak tactics compared to other sources of known jailbreak templates, \ie \danshort \citep{doanytingnow}, \trustllm \citep{sun2024trustllm}, \decodingtrust \citep{wang2023decodingtrust}. \itwshort attacks are also more adversarial than attacks generated by existing semantic-level jailbreak methods (\ie \pair, \tap, \pap) as they, on average, contain more jailbreak tactics per query \citep{pair, tap, pap}. 
Finally, given the diversity of \itwshort jailbreak tactics, it's concerning that existing public safety training data, namely \hhrlhf \citep{hh-rlhf}, \safellamashort \citep{safetytuned_llama}, and \saferlhf \citep{safe-rlhf}, does not contain adversarial enough training examples, limiting downstream models' robustness against adversarial threats. 

